\medskip

\begin{center}\textbf{Partie A : étude d'une fonction} \end{center}

On considère la fonction $f$ définie sur $\R$ par 
\[f(x) = x - \ln \left(x^2 + 1\right),\]

où ln désigne la fonction logarithme népérien.

\medskip

\begin{enumerate}
\item On admet que $f$ est dérivable sur $\R$ et on note $f'$ sa fonction dérivée.
	\begin{enumerate}
		\item Montrer que pour tout nombre réel $x$, on a :
\[f'(x) = \dfrac{(x - 1)^2}{ x^2 + 1}.\]

		\item En déduire le sens de variation de la fonction $f$ sur $\R$.
	\end{enumerate}
\item Montrer que pour tout nombre réel $x > 0$, on a :
\[f(x) = x - 2\ln (x) - \ln \left(1 + \dfrac{1}{x^2}\right).\]

\item Calculer la limite de la fonction $f$ en $+\infty$.
\end{enumerate}

\begin{center}\textbf{Partie B : étude d'une suite}
 \end{center}

On considère la suite $\left(u_n\right)$ définie par : 
\[\left\{\begin{array}{r c l}
u_0& =&7\\
u_{n+1} = f\left(u_n\right)& =&u_n - \ln\left(u_n^2 + 1\right)\:\text{pour tout }\:n \in \N
\end{array}\right.\]

\begin{enumerate}
\item Montrer, en utilisant un raisonnement par récurrence, que pour tout entier naturel $n$ :
$u_n \geqslant 0$.
\item Montrer que la suite $\left(u_n\right)$ est décroissante.
\item En déduire la convergence de la suite $\left(u_n\right)$.
\item On note $\ell$ la limite de la suite $\left(u_n\right)$. Déterminer la valeur de $\ell$. \item
	\begin{enumerate}
		\item Recopier et compléter le script ci-dessous écrit en langage Python afin qu'il
renvoie la plus petite valeur de l'entier $n$ à partir de laquelle $u_n \leqslant h$, où $h$ est un nombre réel strictement positif.

\begin{center}
\begin{tabular}{|l|}\hline
from math import log as ln\\
\#permet d'utiliser la fonction ln \\
\#Le Logarithme népérien\\
~\\
def seuil(h):\\
\quad n = 0\\
\quad u =7\\
\quad while \ldots :\\
\qquad n = n+1\\
\qquad u = \ldots\\
\quad return n\\ \hline
\end{tabular}
\end{center}

		\item Déterminer la valeur renvoyée lorsqu'on saisit seuil(0.01) dans la console Python. Justifier la réponse.
		\end{enumerate}
\end{enumerate}

\begin{center}\textbf{Partie C : calcul intégral}\end{center}

\begin{enumerate}
\item Étudier le signe de la fonction $f$ sur $[0~;~+\infty[$.
\item Interpréter graphiquement l'intégrale:

\[I = \displaystyle\int_2^4 f(x)\:\text{d}x.\]

\item On admet dans cette question que, pour tout nombre réel $x \in [2~;~4]$, on a
l'encadrement:

\[0,5x -1 \leqslant f(x) \leqslant 0,25x + 0,25.\]

En déduire l'encadrement :

\[ 1 \leqslant I \leqslant 2.\]
\end{enumerate}

\bigskip
